\section{Problem Formulation}
\label{sec:formulation}

We study when to compile GUI procedures that agents execute repeatedly while limiting cumulative cost relative to continued agent execution.
Tasks in one family share a GUI procedure but differ in their inputs, e.g., event titles and dates.

We consider two ways to execute, or \emph{serve}, a task.
In \emph{agent execution}, the model selects GUI actions step by step, with mean cost $c$.
In \emph{program execution}, a model call extracts the parameter values from the task prompt, i.e., a \emph{binding}.
The program then executes the shared procedure with these values, without further model calls.
Its model cost $c^{\mathrm{prog}}$ is the extraction cost.
We use weighted tokens: one unit corresponds to one uncached input token at the model's reference price, with cached-input and output tokens weighted by their relative prices.
Appendix~\ref{app:accounting} explains how we calculate these costs.

A compilation attempt uses $k$ completed agent traces.
The attempt charge includes all model calls made during compilation.
Let $p$ be the probability that compilation succeeds, $C$ the mean cost of a successful attempt, and $C^{\mathrm{fail}}$ the mean cost of a failed attempt.
A failed attempt still consumes tokens.
In the online model, these cost values become fixed charges for the corresponding outcomes.
Historical means alone do not bound future actual costs.

\looseness=-1
On each program execution, GUI drift makes the program stop working with probability $h$.
If there is no GUI drift, execution fails with probability $q$.
Either event requires one agent fallback of cost $c$, but only GUI drift stops further program reuse.
The total fallback probability is $h+(1-h)q$, reducing to $q$ when $h=0$.
The cost model assumes detected failures and reports task success separately.
Routing selects an available program or the agent for the current task.
It costs $\tau_0+m\ell_t$ when $\ell_t$ programs are listed, where $\tau_0$ is the cost with an empty list and $m$ is the cost per entry.

\paragraph{Objective.}
We compare PACE with serving every arrival through agent execution, without compiling or reusing programs.
For arrival $t$, define this agent execution reference charge as $a_t=\tau_0+c$ using that family's given cost, and let
\begin{equation}
A_t=\sum_{i=1}^{t}a_i.
\label{eq:agent-reference}
\end{equation}
This reference includes program-served arrivals.
It is a model-based comparison, not an extra agent execution observed on each task.
Let $K_t^\pi$ include policy $\pi$'s routing, serving, fallback, and all compilation charges through arrival $t$, including compilation after the final service.
For an unknown stream length $n$, we seek
\begin{equation}
\min_{\pi\in\Pi_{\mathrm{online}}}\mathbb{E}[K_n^\pi]
\quad\text{subject to}\quad
K_t^\pi\leq(1+\epsilon)A_t,\quad t=1,\ldots,n.
\label{eq:objective}
\end{equation}
The constraint holds for every history allowed by the cost model, while the expectation uses the arrival and outcome model evaluated in the experiments.
The set $\Pi_{\mathrm{online}}$ contains policies that serve each arrival through an agent or an available program, using only the current task, given parameters, and past observations.
We give the policy family labels, cost profiles, routing charges, and $q,h$.
The policy estimates future arrivals of each task family and the probability of successful compilation online.
Looking up a program or removing its list entry costs no tokens.
Savings can fund compilation across families, but future tasks that have not arrived add no budget.
The margin is $\epsilon=0.25$ in the main evaluation.
This objective constrains cost, without imposing a task-success guarantee.
